% Første linje er dokumentklassen. Brug *altid* memoir!
\documentclass[a4paper,oneside,article]{memoir}

% Herunder indlæses forskellige pakker
\usepackage[utf8]{inputenc}  % Korrekt håndtering af æ, ø og å
\usepackage[T1]{fontenc}  % Korrekt håndtering af æ, ø og å
\usepackage{microtype}  % Typografisk magi! Giver bl.a. pænere orddeling
\usepackage{graphicx}  % Gør det muligt at indsætte billeder
\usepackage{amsmath}  % Giver adgang til uundværlige matematikting
\usepackage{mathtools}  % Retter ting i ams og giver ekstra muligheder
\usepackage[danish]{babel}  % Danske betegnelser og orddeling
\usepackage{lmodern}
\usepackage{alphabeta}
\usepackage{float} %Mere kontrol over billede-placering
\renewcommand{\danishhyphenmins}{22}  % Bedre dansk orddeling
\usepackage[exponent-product=\ensuremath{\cdot}]{siunitx} % Sientific notation for forskellige ting, med en regel om at den skal bruge cdot istedet for times. 
\usepackage{derivative} % Til at skrive differentialligninger 
\usepackage{xcolor} % For at få flere tekst farver 
\usepackage[colorlinks=true]{hyperref} % For at få refrences man kan trykke på, slet [] hvis man vil have kasser i stedet for farvet tekst.
\usepackage{pdfpages} % Til at indsætte andre pdf'er
\usepackage{subfig} % For at have to billeder ved siden af hinanden
\usepackage{unicode-math}

% Pænere toc indents 
\usepackage{tocloft}
\cftsetindents{section}{1em}{2em}
\cftsetindents{subsection}{4em}{0em}


% Indholdet i dit dokument starter her!
\begin{document}
\author{Mikkel Moth Billing \\ 202307989}
\title{Machine Learning}
\date{\today}
\maketitle   % Indsætter overskriften
% Table of Contents indstillinger
\setcounter{tocdepth}{2}
\setcounter{secnumdepth}{1}
\hypersetup{linkcolor=black} % Laver tableofcontents sort.  
\tableofcontents %Indsætter ToC
\hypersetup{linkcolor=red} % og så alle andre hyperrefs røde. 

%Custom Commands
%Til opgavebeskrivelser
\newcommand{\opg}[1]{\textcolor{gray}{\emph{#1}}\plainbreak{1}} 

\newcommand{\ti}[1]{\cdot 10^{#1}} % Let: gange 10^x
\newcommand{\e}[1]{\emph{e}^{#1}} % Pænere e^x
\newcommand{\vm}[2]{\begin{pmatrix} #1 \\ #2 \end{pmatrix}}
\newcommand{\mat}[1]{\begin{pmatrix} #1 \end{pmatrix}}

% Lette paranteser 
\newcommand{\br}[2][1]{%
    \ifcase#1
    \or \left( #2 \right) %1 eller ingenting
    \or \left[ #2 \right] %2
    \or \left\langle #2 \right\rangle %3
    \or \left\{ #2 \right\} %4
    \or \left| #2 \right| %5
    \fi
}
%Til at sætter bileder ind i store dokumenter
\newcommand{\bil}[1]{
\begin{figure}[H]
    \centering
    \includegraphics[width=0.75\linewidth]{Billeder/#1.png}
    \label{fig:#1}
\end{figure}
}

%Til at sætte pdf'er ind i store dokumenter (kun pdf'er på mere end 1 sider)
\newcommand{\pdf}[4]{
\includepdf[pages=#3,scale=.85,pagecommand={\subsection{#2 - \ref{ch:opg}}  \label{asc:#2}}]{pdf/#1.pdf}
\includepdf[pages=\the\numexpr 1 + #3\relax-#4, scale=.85]{pdf/#1.pdf}
}

% Lav ny pagestyle (kun odd fordi enkeltsidet)
\makepagestyle{Title}
\pagestyle{Title}
%\makeoddhead{Title}{\footnotesize\color{gray}{\theauthor}}{}{\textcolor{gray}{\thedate}}
\makeoddfoot{Title}{}{\footnotesize\color{gray}{\thepage{} / \thelastpage}}{}
% Også fod på siden med \maketitle
\makeoddfoot{plain}{}{\footnotesize\color{gray}{\thepage{} / \thelastpage}}{}
\newpage

%Skriv
\chapter{Linear model}
Hypothesis consists of the features vector: 
\begin{align*}
    x = \br[4]{x_1,..., x_d}
\end{align*}
and a bias $b$. 
if the linear combenation of the wieghts and the input, 
and compare it to the bias, we can make predictions. 

Example of mortgage, approve is $1$ and decline is $-1$. 
So we can predict that we approve
\begin{align*}
    w_1x_1+...+w_dx_d > -b 
\end{align*}
\bil{Linear}

To find the parameters we can use \textbf{Perceptron Learning Algorithm (PLA)}.

Algorithm: 
While there is a data point (x,y) with \[y \neq \text{sign}(w^Tx) \rightarrow w = w +yx.\]
\bil{PLA1} 

As we can see the algorithm can enter an infinite loop, 
as there might not be a linearly separable dataset. 

A way to fix this is to use Pocket PLA which saves the most accurate hyperplane. 
\bil{PocketPLA}

\section{Decision Trees}
\bil{dTree}
We would like to keep the decision tree small to not overfit to the training data. 
So we could find the best performing tree with $k$ nodes. 
But this is \textcolor{red}{NP-hard}, so we rely on \textcolor{red}{heureistics}. 
With a greedy learning algorithm ith $k=2$: 
\bil{dtree2}
We test the tree on all possible split, this has $O(n^2d)$ runtime. 

Another way to build a tree is to increase $k$ to maximise the accuracy. 
\bil{dtree3}

\subsection{Entropy}
Another way to test decision tree is to base it on entropy. 
If we have $n$ point with $n_i$ as the numbers of points with label $i$. 
And $p_i = n/n_i$. Then entropy is: 
\begin{align*}
    H(S) = -\sum_{i=1}^{k} p_i\log_2 p_i. 
\end{align*}

Entropy is a measure of uncertainty in distrubtions, high entropy is high unsertainty. 
$H(S)$ describes the expected number of bits it takes to describe the label of a random element drawn from S. 
If all elements in $S$ er equal $H(S) = 0$. If there is an equal amount $H(S) = \log_2k$.
\bil{entropy1}
\bil{entropy2}

\subsection{Loss Function}
A Loss Function weighs how much a mistake influences a model. 
We can then try to minimize the loss function over the training points. 
An example is the Least Square Loss Function: 
\begin{align*}
    L(y',y) = (y'-y)^2.
\end{align*}
\bil{lf1}



\end{document} % ----------Dokumentet er nu slut! ------------